\chapter{Introduction to the Course [10/09/2026]}
The objective of the course is to provide principles, techniques, and tools for Security Testing.

The project, as part of the exam, is written in Java, but the course also explores other programming languages like C/C++, Python, PHP, and SQL. This is because different programming languages suffer from different vulnerabilities.

\section{Contents of the Course}
The main topics of the course are described below:

\paragraph{Introduction to Security Testing} An overview of the security testing life-cycle, vulnerabilities, bugs, and secure coding.
\paragraph{Context} A context regarding vulnerability analysis, rates, and repositories will be provided.
\paragraph{Vulnerabilities} A detailed presentation of SQL Injection, XSS, String formatting issues, and more. 
\paragraph{Vulnerability Detection} Both static and dynamic analysis will be discussed, and an overview of AI-based vulnerability detection will be provided.
\paragraph{Tools} For these analyses, tools need to be provided; they will be analyzed and discussed.
\paragraph{[OPTIONAL] Advanced Topics} 

\section{Exam}
The exam is composed of two parts:
\begin{itemize}
    \item A written test with exercises as well as open and closed-form questions (maximum two (2) hours) - About 60\% of the grade.
    \item A written project report \textbf{to be submitted before the written exam} - About 40\% of the grade.
\end{itemize}
It is permitted to attend both in one session or in distinct sessions with no expiration date. Both parts need to be at least 18/30; grades are expressed out of 30.

A simulation of the written part will be performed. Some laboratories will be devoted to working on the exam project.

\subsection{Project Rules}
\begin{itemize}
    \item Groups can be composed of one or two students.
    \item The project can be submitted a maximum of two non-consecutive times.
    \item If the project is graded as a fail two times, a different test has to be arranged.
    \item A passing project can be used for a maximum of two written tests.
    \item In case of a passing final grade, the project cannot be resubmitted.
    \item In case of resubmission, a clear difference/delta between the first and the second submission must exist; otherwise, the final score will be reduced by 1.
\end{itemize}

\textbf{TOTAL OR PARTIAL PLAGIARISM IS FORBIDDEN}

\noindent This applies to both the project and the written test.